%% PLN credal degrees of freedom
%%
%% Render with:
%%   lualatex -shell-escape pln-credal-degrees-of-freedom.tex \
%%     && bibtex pln-credal-degrees-of-freedom \
%%     && lualatex -shell-escape pln-credal-degrees-of-freedom.tex \
%%     && lualatex -shell-escape pln-credal-degrees-of-freedom.tex

\documentclass[]{article}
\usepackage{url}
\usepackage[hyperindex,breaklinks]{hyperref}
\IfFileExists{breakurl.sty}{\usepackage{breakurl}}{}
\usepackage{listings}
\lstset{basicstyle=\ttfamily\footnotesize,breaklines=true,frame=single,
        xleftmargin=1em,xrightmargin=0.5em}
\usepackage{float}
\restylefloat{table}
\usepackage{longtable}
\usepackage{booktabs}
\usepackage{graphicx}
\usepackage[font=small,labelfont=bf]{caption}
\usepackage[skip=0pt]{subcaption}
\usepackage{amsfonts}
\usepackage{amsmath}
\usepackage{amsthm}
\usepackage{unicode-math}
\IfFileExists{newunicodechar.sty}{
  \usepackage{newunicodechar}
  \newunicodechar{∘}{\ensuremath{\circ}}
}{}
\IfFontExistsTF{Stix Two Math}{\setmathfont{Stix Two Math}}{\setmathfont{STIX Math}}

% Compatibility fallback for older TeX bundles.
\makeatletter
\providecommand{\headerps@out}{}
\makeatother

\newtheorem{theorem}{Theorem}
\newtheorem{lemma}[theorem]{Lemma}
\newtheorem{proposition}[theorem]{Proposition}
\newtheorem{corollary}[theorem]{Corollary}
\newtheorem{definition}[theorem]{Definition}
\newtheorem{remark}[theorem]{Remark}
\newtheorem{observation}[theorem]{Observation}

\newcommand{\nuPLN}{\texorpdfstring{$\nu$}{nu}\textup{PLN}}
\newcommand{\xiPLN}{\texorpdfstring{$\xi$}{xi}\textup{PLN}}
\newcommand{\code}[1]{\texttt{#1}}
\newcommand{\codepath}[1]{\path{#1}}
\newcommand{\lid}[1]{{\ttfamily\nolinkurl{#1}}}
\newcommand{\lpath}[1]{\texttt{\nolinkurl{#1}}}
\newcommand{\ty}[1]{\ensuremath{\mathtt{#1}}}
\newcommand{\STV}[2]{\langle #1,\, #2 \rangle}
\newcommand{\strength}{\mathrm{strength}}
\newcommand{\confidence}{\mathrm{confidence}}
\newcommand{\Eq}{\;\Leftrightarrow\;}

\emergencystretch=3em

\title{Credal degrees of freedom for PLN truth values:\\
       a layered characterization of strength and confidence}
\author{Mettapedia Working Notes}
\date{\today}

\begin{document}
\maketitle

\begin{abstract}
PLN's two-number truth value $\STV{s}{c}$ is the compression of a richer
object: a credal set of admissible precise previsions over an outcome space.
We give a layered formal characterization of where $\STV{s}{c}$ is
\emph{determined} by the spec --- the strength is a single real, the
confidence is at its information ceiling --- and where the underlying
credal set has \emph{degrees of freedom} that the two-number summary
cannot see.

The characterization rests on a recently-landed Lean dichotomy in
\codepath{Mettapedia.ProbabilityTheory.ImpreciseProbability.ProjectiveCredal}:
the predicates \lid{credalSetDetermines} and \lid{credalSetHasStrictWidth}
on a credal set $C$ and a gamble $X$ are the formal locators of, respectively,
``$\STV{s}{c}$ is fully determined on $X$'' and ``$\STV{s}{c}$ has
genuine interval width on $X$.'' The collapse theorem
\lid{lower\_eq\_upperEnvelope\_of\_credalSetDetermines} closes the
``determined $\Rightarrow$ precise'' direction; the width theorem
\lid{lower\_upperEnvelope\_nontrivial\_of\_strictWidth} closes the
``disagreement $\Rightarrow$ envelope opens'' direction.

We lift this dichotomy to the projective credal level
(\lid{determinesGlobalGamble} / \lid{hasStrictGlobalWidth}) and onward to
the infinite Markov-logic setting through DLR completions
(\lid{dlrQueryDetermined} / \lid{dlrQueryHasStrictWidth}). Under the
Dobrushin uniform small-influence hypothesis,
\lid{PaperUniformSmallTotalInfluence} implies \lid{dlrQueryDetermined},
and therefore the MLN query envelope is precise. In the opposite
direction, whenever two DLR completions are explicitly exhibited that
disagree on a finite query, the corresponding Walley envelope has strict
width. Failure of the Dobrushin hypothesis is therefore a frontier where
strict width can occur, not by itself a proof that it must occur.

We organize the degrees of freedom into three slots:
(A)~specification-level freedom in the local credal sets;
(B)~completion-level freedom in extreme-point choice within a fixed credal
set (aleatoric, formally invisible to $\STV{s}{c}$);
(C)~phase-transition freedom from non-uniqueness of DLR completions
in infinite MLNs. We honestly enumerate what is open.
\end{abstract}

\tableofcontents

\section{Introduction}

PLN's atomic truth-value object is a pair $\STV{s}{c}$ with $s, c \in [0,1]$
called \emph{strength} and \emph{confidence}. Standard PLN parametrises
$s$ as a probability estimate and $c$ as a function of evidence count,
typically $c = N/(N+\kappa)$ for a system-level constant $\kappa$ and an
evidence-count integer $N$. The pair is what travels through PLN's
rules (revision, deduction, induction, abduction); it is what is read off
by query atoms; it is the type signature of every PLN belief.

But $\STV{s}{c}$ is a two-number summary. The same pair can be produced by
very different underlying epistemic objects: a Beta-distribution
second-order Bayesian; a Walley-style coherent lower prevision with a
specific interval width; an exchangeable de Finetti mixture with a
particular concentration. PLN's compositional rules treat $\STV{s}{c}$
\emph{as} the belief object, and the algebra reduces to numerical update
formulas; but the question of \emph{what is and is not pinned down by
$\STV{s}{c}$} is rarely written explicitly in the PLN literature.

This paper makes that explicit. We organise the question as follows.
The strength and confidence of a query in PLN are determined by an
underlying credal set $C$ of admissible precise previsions. The
\emph{degrees of freedom} that $\STV{s}{c}$ summarises but does not
expose live in three slots:
\begin{itemize}
\item Specification freedom in the local credal sets that build $C$;
\item Completion freedom (extreme-point choice) inside a fixed $C$;
\item Phase-transition freedom from non-uniqueness of global completions
      when $C$ is the projective limit of a Markov-logic / DLR family.
\end{itemize}
Each slot has a formally locatable status. Recent landings in
the \codepath{Mettapedia.ProbabilityTheory.ImpreciseProbability} and
\codepath{Mettapedia.Logic} libraries (the \emph{credal dichotomy}
predicates and the DLR query lifts; see Sec.~\ref{sec:dichotomy} and
\ref{sec:dlr}) make this dichotomy an actual Lean theorem rather than
informal prose. Wherever possible we quote the Lean theorem statements
verbatim from their source files.

\paragraph{Scope.}
We treat strength and confidence as the read-off of credal-envelope width
and centre. We do \emph{not} address the propagation of $\STV{s}{c}$
through PLN's compositional rules (deduction, induction, abduction,
revision); that material is in the existing PLN corpus
(\xiPLN, \nuPLN, the WM-PLN book; see Sec.~\ref{sec:related}). The paper's
contribution is the determinacy / freedom characterization at the level
of a single truth value, with formal Lean hooks.

\paragraph{Sibling papers.}
This paper is a sibling of \xiPLN\ and the \emph{PLN Review}
(\lpath{papers/xiPLN.tex}, \lpath{papers/pln-review.tex}). The
WM-PLN book \lpath{papers/wm-pln-book.tex} contains an informal section
on degrees of freedom versus forcing (\S 7.4 in the current draft); the
present paper offers the formal characterization that the book section
calls for.

\section{The credal stack: types and definitions}\label{sec:stack}

We work in the formal stack of
\codepath{Mettapedia.ProbabilityTheory.ImpreciseProbability.ProjectiveCredal}
(referenced throughout as \emph{ProjectiveCredal.lean}, file path
\lpath{Mettapedia/ProbabilityTheory/ImpreciseProbability/ProjectiveCredal.lean}).
Outcome spaces are arbitrary types $\Omega$; gambles are bounded real-valued
functions on $\Omega$.

\subsection{Gambles and previsions}

A \emph{gamble} on $\Omega$ is a real-valued function on $\Omega$:
\begin{lstlisting}
abbrev Gamble (Omega : Type*) := Omega -> Real
\end{lstlisting}

A \emph{precise prevision} on $\Omega$ is a Walley-coherent linear
functional from gambles to reals: the structure
\lid{PrecisePrevision} in ProjectiveCredal.lean is the standard linear
expectation with constant-preservation, positive homogeneity, and
additivity built in:
\begin{lstlisting}
structure PrecisePrevision (Omega : Type*) where
  toFun       : Gamble Omega -> Real
  lower_bound : ... -- c <= X (pointwise) implies c <= toFun X
  pos_homog   : ... -- toFun (r . X) = r * toFun X for r >= 0
  add         : ... -- toFun (X + Y) = toFun X + toFun Y
\end{lstlisting}

A \emph{lower prevision} relaxes additivity to superadditivity, and is
the Walley credal-summary object:
\begin{lstlisting}
structure LowerPrevision (Omega : Type*) where
  toFun       : Gamble Omega -> Real
  lower_bound : ... -- c <= X pointwise implies c <= toFun X
  pos_homog   : ... -- toFun (r . X) = r * toFun X for r >= 0
  superadd    : ... -- toFun (X + Y) >= toFun X + toFun Y
\end{lstlisting}
Coherence (positive homogeneity, superadditivity, and avoidance of sure
loss) is the Walley condition for the lower prevision to be the
lower envelope of some non-empty set of precise previsions.

\subsection{Credal sets and envelopes}

A \emph{credal set} is just a set of precise previsions:
\begin{lstlisting}
abbrev CredalPrevisionSet (Omega : Type*) := Set (PrecisePrevision Omega)
\end{lstlisting}

For a credal set $C$ and a gamble $X$, the \emph{lower} and \emph{upper}
envelopes are the standard inf / sup:
\begin{lstlisting}
def lowerEnvelope (C : CredalPrevisionSet Omega) (X : Gamble Omega) : Real :=
  sInf ((fun P : PrecisePrevision Omega => P X) '' C)

def upperEnvelope (C : CredalPrevisionSet Omega) (X : Gamble Omega) : Real :=
  sSup ((fun P : PrecisePrevision Omega => P X) '' C)
\end{lstlisting}
The construction
\lid{lowerEnvelopePrevision} packages \lid{lowerEnvelope} together with
the proofs of constant-preservation, positive homogeneity, and
superadditivity, exhibiting a \lid{LowerPrevision} object out of any
non-empty bounded-below credal set. The finite case theorem
\lid{lowerEnvelopePrevision\_isCoherent\_of\_finite}
(ProjectiveCredal.lean line~868) confirms Walley coherence on finite
nonempty $\Omega$.

\subsection{PLN $\STV{s}{c}$ as the credal-envelope summary}\label{sec:stv-summary}

The canonical Lean STV type lives in
\codepath{Mettapedia.Logic.PLNDeduction} at lines 138--139:
\begin{lstlisting}
/-- Simple Truth Value: (strength, confidence) pair -/
structure STV where
  strength   : Real
  confidence : Real
  strength_nonneg  : 0 <= strength
  strength_le_one  : strength <= 1
  confidence_nonneg : 0 <= confidence
  confidence_le_one : confidence <= 1
\end{lstlisting}
Its distributional model is in \codepath{Mettapedia.Logic.PLNDistributional}:
given a Beta$(\alpha, \beta)$ distribution, the canonical PLN STV reads
\begin{equation*}
   s \;=\; \frac{\alpha}{\alpha + \beta},
   \qquad
   c \;=\; f(\alpha + \beta)
\end{equation*}
for an evidence-count map $f$ (typically $f(N) = N/(N+\kappa)$).

The two-number compression from a credal set $C$ and a query gamble
$X_q$ (e.g.\ the indicator $\mathbf{1}_A$ of a query atom $A$) is the
following implicit map:
\begin{equation}\label{eq:summary}
   C, X_q \;\longmapsto\;
   \left\{
   \begin{array}{l}
   s_q \;=\; \tfrac{1}{2}\bigl(\,\mathrm{lowerEnvelope}\, C\, X_q\,
                              +\,\mathrm{upperEnvelope}\, C\, X_q\,\bigr) \\[4pt]
   c_q \;=\; g\bigl(\,\mathrm{upperEnvelope}\, C\, X_q\,
                   -\,\mathrm{lowerEnvelope}\, C\, X_q\,\bigr)
   \end{array}
   \right.
\end{equation}
for a monotone-decreasing $g : [0,1] \to [0,1]$ with $g(0) = 1$ (no
width, full confidence) and $g(1) = 0$ (maximum width, no information).
The specific choice of midpoint and $g$ is an interface design decision
(PLN traditionally uses midpoint for $s$ and a sigmoidal $g$ derived from
the Beta variance); the freedom characterization of this paper is
independent of this choice. What matters is that \eqref{eq:summary}
is the read-off, and that the credal set $C$ may have degrees of
freedom that the pair $(s_q, c_q)$ cannot expose.

\begin{remark}
The summary map \eqref{eq:summary} loses two kinds of information.
First, \emph{which} gambles have positive credal-width and which do
not: $(s_q, c_q)$ summarises a single query but a credal set with
non-zero width on \emph{some} queries may have zero width on others.
Second, the \emph{shape} of $C$: a credal set on $\Omega$ whose extreme
points are clustered (high agreement) and one whose extreme points are
dispersed (low agreement) can produce the same midpoint and width on a
particular query, but compose differently under PLN's inference rules.
The width as a function on gambles --- which we will study below --- is
the right object; $(s_q, c_q)$ is its evaluation at one point.
\end{remark}

\section{Finite-case determinacy}\label{sec:finite}

We first treat the case where the outcome space $\Omega$ is finite. Here
the Walley theory is complete in the formalization; we summarise the
load-bearing results.

\subsection{Uniform positivity and avoidance of sure loss}

For finite $\Omega$, a pointwise-strictly-positive gamble has a uniform
positive lower bound (ProjectiveCredal.lean line~833):
\begin{lstlisting}
theorem finite_strictlyPositive_uniformLowerBound
    {Omega : Type*} [Fintype Omega] [Nonempty Omega]
    (X : Gamble Omega) (hX : forall omega, 0 < X omega) :
    exists epsilon : Real, 0 < epsilon and forall omega, epsilon <= X omega
\end{lstlisting}
This is the technical hinge that lets every nonempty finite credal set
\emph{avoid sure loss} in the strict Walley sense
(line~855):
\begin{lstlisting}
theorem lowerEnvelopePrevision_avoidsSureLoss_of_finite
    {Omega : Type*} [Fintype Omega] [Nonempty Omega]
    (C : CredalPrevisionSet Omega) (hC : C.Nonempty)
    (hBdd : forall X : Gamble Omega,
      BddBelow ((fun P : PrecisePrevision Omega => P X) '' C)) :
    (lowerEnvelopePrevision C hC hBdd).avoidsSureLoss
\end{lstlisting}
and to be \emph{coherent} (line~868):
\begin{lstlisting}
theorem lowerEnvelopePrevision_isCoherent_of_finite
    {Omega : Type*} [Fintype Omega] [Nonempty Omega]
    (C : CredalPrevisionSet Omega) (hC : C.Nonempty)
    (hBdd : forall X : Gamble Omega,
      BddBelow ((fun P : PrecisePrevision Omega => P X) '' C)) :
    (lowerEnvelopePrevision C hC hBdd).isCoherent
\end{lstlisting}
The take-away for PLN: on a finite outcome space, the credal-envelope
construction never fails. The map from spec to $\STV{s}{c}$ is total.

\subsection{Singleton collapse}

The case where the credal set has exactly one element is the case where
PLN's $\STV{s}{c}$ has no freedom at all; everything is pinned. The
theorem (ProjectiveCredal.lean line~885) is direct:
\begin{lstlisting}
theorem lowerEnvelope_singleton
    {Omega : Type*} (P : PrecisePrevision Omega) (X : Gamble Omega) :
    lowerEnvelope ({P} : CredalPrevisionSet Omega) X = P X
\end{lstlisting}
and similarly for \lid{upperEnvelope\_singleton}. On a singleton credal
set, lower = upper on every gamble. PLN-side: every query has $c$ at its
information ceiling, and $s$ is the single value $P X_q$.

\section{The credal dichotomy: where strength is determined and where it has width}\label{sec:dichotomy}

The contribution of the recent landing in ProjectiveCredal.lean is to
make the ``determined / has width'' distinction \emph{predicate-level},
not theorem-level. Two definitions are added (lines~746 and~753), each
followed by collapse / non-triviality theorems that pin down the
PLN-side consequence.

\subsection{The two predicates}

\begin{definition}[determined credal set on a gamble]\label{def:determines}
The predicate \lid{credalSetDetermines} (ProjectiveCredal.lean line~746):
\begin{lstlisting}
def credalSetDetermines {Omega : Type*}
    (C : CredalPrevisionSet Omega) (X : Gamble Omega) : Prop :=
  forall P : PrecisePrevision Omega, P in C ->
    forall Q : PrecisePrevision Omega, Q in C -> P X = Q X
\end{lstlisting}
$C$ \emph{determines} $X$ when every two admissible precise previsions in
$C$ assign the same expectation to $X$.
\end{definition}

\begin{definition}[strict width]\label{def:strictwidth}
The predicate \lid{credalSetHasStrictWidth} (line~753):
\begin{lstlisting}
def credalSetHasStrictWidth {Omega : Type*}
    (C : CredalPrevisionSet Omega) (X : Gamble Omega) : Prop :=
  exists P : PrecisePrevision Omega, P in C and
    exists Q : PrecisePrevision Omega, Q in C and P X < Q X
\end{lstlisting}
$C$ has \emph{strict width} on $X$ when there exist two precise previsions
in $C$ that strictly disagree on $X$.
\end{definition}

The two predicates are mutually exclusive (ProjectiveCredal.lean line~809):
\begin{theorem}[mutual exclusion]\label{thm:mutex}
\lid{not\_credalSetDetermines\_of\_strictWidth}:
\begin{lstlisting}
theorem not_credalSetDetermines_of_strictWidth
    {Omega : Type*}
    {C : CredalPrevisionSet Omega} {X : Gamble Omega}
    (hWidth : credalSetHasStrictWidth C X) :
    not (credalSetDetermines C X)
\end{lstlisting}
\end{theorem}

The endpoint instances cover the trivial cases. Subsingleton credal sets
collapse to determinacy
(\lid{credalSetDetermines\_of\_subsingleton}, line~758), and so do
singletons (\lid{credalSetDetermines\_singleton}, line~766):
\begin{lstlisting}
theorem credalSetDetermines_singleton {Omega : Type*}
    (P : PrecisePrevision Omega) (X : Gamble Omega) :
    credalSetDetermines ({P} : CredalPrevisionSet Omega) X
\end{lstlisting}

\subsection{Collapse and width theorems}

The PLN-side consequence of \lid{credalSetDetermines} is precisely
that strength is pinned: the lower and upper envelopes coincide
(ProjectiveCredal.lean line~798):
\begin{theorem}[envelope collapse from determinacy]\label{thm:collapse}
\lid{lower\_eq\_upperEnvelope\_of\_credalSetDetermines}:
\begin{lstlisting}
theorem lower_eq_upperEnvelope_of_credalSetDetermines
    {Omega : Type*}
    (C : CredalPrevisionSet Omega) (X : Gamble Omega)
    (hC : C.Nonempty)
    (hBddBelow : BddBelow ((fun P => P X) '' C))
    (hBddAbove : BddAbove ((fun P => P X) '' C))
    {P : PrecisePrevision Omega} (hP : P in C)
    (hDet : credalSetDetermines C X) :
    lowerEnvelope C X = upperEnvelope C X
\end{lstlisting}
This is the formal statement of ``$\STV{s}{c}$ on $X_q$ is precise''
when $C$ determines $X_q$.
\end{theorem}

The PLN-side consequence of \lid{credalSetHasStrictWidth} is the opposite:
the envelope is genuinely non-trivial (line~835):
\begin{theorem}[envelope width from strict disagreement]\label{thm:width}
\lid{lower\_upperEnvelope\_nontrivial\_of\_strictWidth}:
\begin{lstlisting}
theorem lower_upperEnvelope_nontrivial_of_strictWidth
    {Omega : Type*}
    (C : CredalPrevisionSet Omega) (X : Gamble Omega)
    (hBddBelow : BddBelow ((fun P => P X) '' C))
    (hBddAbove : BddAbove ((fun P => P X) '' C))
    (hWidth : credalSetHasStrictWidth C X) :
    lowerEnvelope C X < upperEnvelope C X
\end{lstlisting}
This is the formal statement of ``$\STV{s}{c}$ on $X_q$ is genuinely an
interval'' when two completions in $C$ strictly disagree on $X_q$.
\end{theorem}

\subsection{PLN read-off}\label{sec:plnreadoff}

Putting the dichotomy into the summary map \eqref{eq:summary}: for any
credal set $C$ on which both envelope-boundedness conditions hold, and any
query gamble $X_q$, exactly one of the following holds:
\begin{center}
\begin{tabular}{lll}
\toprule
predicate & envelope & PLN $\STV{s}{c}$ \\
\midrule
\lid{credalSetDetermines} $C$ $X_q$ & $L = U$ & precise $s = L\,X_q$; $c$ at ceiling \\
\lid{credalSetHasStrictWidth} $C$ $X_q$ & $L < U$ & $s$ in interval $[L, U]$; $c < $ ceiling \\
\bottomrule
\end{tabular}
\end{center}
The Lean dichotomy is the formal locator of which case the PLN truth-value
read-off lives in.

\section{Projective credal specs: three slots for freedom}\label{sec:projective}

Once $\Omega$ is large (or infinite), one builds the credal set
\emph{compositionally} from finite-window data. The standard packaging
is a projective family of local credal sets, indexed by ``windows''
(finite subsets, finite cylinders) of the global outcome space, glued
by restriction maps.

\subsection{The projective machinery}

The cylinder structure (ProjectiveCredal.lean, the
\lid{ProjectiveCylinderSystem} structure) packs the window-to-local
projection and the local restriction maps:
\begin{lstlisting}
structure ProjectiveCylinderSystem (Window Global : Type*) [LE Window] where
  Local             : Window -> Type*
  project           : forall i : Window, Global -> Local i
  restrict          : forall {i j : Window}, i <= j -> Local j -> Local i
  project_restrict  : ...   -- compatibility on the diagram
\end{lstlisting}
The full spec layers a local credal choice on top:
\begin{lstlisting}
structure ProjectiveLocalCredalSpec (Window Global : Type*) [LE Window] where
  cylinders   : ProjectiveCylinderSystem Window Global
  localCredal : forall i : Window, CredalPrevisionSet (cylinders.Local i)
\end{lstlisting}
The \emph{global credal set} associated to a spec is the
\lid{projectiveLimitCredalSet} (the set of $P$ in \lid{PrecisePrevision}
\lid{Global} that marginalise into \lid{localCredal}{} $i$ at every window
$i$).

\subsection{Three identifying cases and the dichotomy at the global level}

The codebase ships two endpoints and one general construction:
\begin{itemize}
\item \textbf{Singleton} (ProjectiveCredal.lean line~1615):
\begin{lstlisting}
def singletonIdentityProjectiveSpec {Omega : Type*}
    (P : PrecisePrevision Omega) : ProjectiveLocalCredalSpec PUnit Omega
\end{lstlisting}
Every window's local credal is the supplied singleton $\{P\}$. The
projective limit is $\{P\}$
(\lid{singletonIdentityProjectiveSpec\_projectiveLimitCredalSet},
line~1620). PLN $\STV{s}{c}$ on every gamble is precise; this is the
``no freedom'' case.

\item \textbf{Unrestricted} (line~1571):
\begin{lstlisting}
def finiteUnrestrictedProjectiveSpec (Omega : Type*) :
    ProjectiveLocalCredalSpec PUnit Omega
\end{lstlisting}
Every window's local credal is \lid{Set.univ}. The projective limit
contains every precise prevision on $\Omega$. PLN $\STV{s}{c}$ on most
gambles has \lid{credalSetHasStrictWidth}; this is the ``maximum freedom''
case. The non-emptiness of the projective limit is witnessed by
\lid{finiteUnrestricted\_finiteEvaluationCompact\_hasCompatibleCompletion}
(line~1578) via a compact carrier in the finite-evaluation topology.

\item \textbf{Genuine intermediate}: any spec whose local credals are
non-trivial proper subsets of \lid{Set.univ}. Whether
\lid{credalSetDetermines} or \lid{credalSetHasStrictWidth} holds on a
particular gamble depends on the spec and the gamble.
\end{itemize}

The dichotomy lifts to the projective level. ProjectiveCredal.lean defines
\lid{determinesGlobalGamble} (line~1307) and \lid{hasStrictGlobalWidth}
(line~1314) as the global analogues of \lid{credalSetDetermines} and
\lid{credalSetHasStrictWidth}, with the mutual exclusion theorem
\lid{not\_determinesGlobalGamble\_of\_strictWidth} at line~1411. The
natural-extension lifting theorems (around lines~1391 and~1424) connect
the projective predicates to the lower / upper envelopes on the projective
limit.

\subsection{The three formal slots, structurally}\label{sec:slots-structural}

We can now name the three slots in which PLN credal freedom lives. Each
is formally locatable through one of the predicates above.

\begin{description}
\item[(A) Specification-level freedom.] The choice of \lid{localCredal}{} $i$
at each window. This is the modeller's freedom: evidence width, prior
choices, and what classes of measure are taken to be admissible at each
local scale.

\item[(B) Completion-level freedom.] Given a fixed credal set $C$, the
choice of $P \in C$ that actually governs the world. This freedom is
purely aleatoric: it is invisible to $\STV{s}{c}$ because $\STV{s}{c}$ is
a property of $C$, not of any specific $P \in C$.

\item[(C) Phase-transition freedom.] For globally-glued projective specs,
the projective limit may itself have multiple admissible elements that
correspond to different \emph{global} solutions of the same local
problem. This is the slot Markov-logic / DLR exemplifies (next section).
\end{description}

Slots (A) and (C) are the only ones $\STV{s}{c}$ can see. Slot (B) is
formally invisible at the truth-value level.

\section{Infinite Markov logic: Dobrushin collapse and DLR width}\label{sec:dlr}

For the infinite case --- countably-many ground atoms, Markov-logic
network with weights, density specified locally --- the projective
machinery specialises into the Dobrushin--Lanford--Ruelle framework. The
file
\codepath{Mettapedia.Logic.MarkovLogicInfiniteCredalBridge}
(\lpath{Mettapedia/Logic/MarkovLogicInfiniteCredalBridge.lean}) packages this.

\subsection{Classical MLN spec and DLR completions}

A \lid{ClassicalInfiniteGroundMLNSpec} packages a countable atom set, a
clause-id set, the clauses, their log-weights, and the local
finite-volume region structure. A \emph{DLR completion} of such a spec
is a probability measure on the infinite world satisfying the
DLR consistency condition at every finite volume:
\begin{lstlisting}
abbrev DLRCompletion (M : ClassicalInfiniteGroundMLNSpec Atom ClauseId) :=
  { mu : ProbabilityMeasure (InfiniteWorld Atom)
  // FixedRegionCylinderDLR M.toStrictlyPositiveInfiniteGroundMLNSpec
       (mu : Measure (InfiniteWorld Atom)) }
\end{lstlisting}
The set of DLR completions of $M$ is the ``credal set'' for the infinite
spec, and the projective-credal lift packages it as
\lid{DLRProjectiveCredalSpecialization} in
\codepath{Mettapedia/Logic/MarkovLogicInfiniteCredalBridge.lean}.

A \emph{query} is a finite-volume event; its \emph{probability} under a
DLR completion is its measure:
\begin{lstlisting}
noncomputable def dlrCompletionQueryProb
    (M : ClassicalInfiniteGroundMLNSpec Atom ClauseId)
    (q : ConstraintQuery Atom) (mu : DLRCompletion M) : Real :=
  ENNReal.toReal
    ((infiniteMLNMassSemantics M mu.1 mu.2).queryProb q)
\end{lstlisting}

\subsection{DLR-level dichotomy}\label{sec:dlrdichotomy}

The credal dichotomy lifts directly to the DLR setting:
\begin{lstlisting}
def dlrQueryDetermined
    (M : ClassicalInfiniteGroundMLNSpec Atom ClauseId)
    (q : ConstraintQuery Atom) : Prop :=
  forall mu nu : DLRCompletion M,
    dlrCompletionQueryProb M q mu = dlrCompletionQueryProb M q nu

def dlrQueryHasStrictWidth
    (M : ClassicalInfiniteGroundMLNSpec Atom ClauseId)
    (q : ConstraintQuery Atom) : Prop :=
  exists mu : DLRCompletion M, exists nu : DLRCompletion M,
    dlrCompletionQueryProb M q mu < dlrCompletionQueryProb M q nu
\end{lstlisting}
\lid{dlrQueryDetermined} says all DLR completions agree on the query;
\lid{dlrQueryHasStrictWidth} says two of them strictly disagree. The
mutual exclusion theorem
\lid{not\_dlrQueryDetermined\_of\_strictWidth}:
\begin{lstlisting}
theorem not_dlrQueryDetermined_of_strictWidth
    (M : ClassicalInfiniteGroundMLNSpec Atom ClauseId)
    (q : ConstraintQuery Atom)
    (hWidth : dlrQueryHasStrictWidth M q) :
    not (dlrQueryDetermined M q)
\end{lstlisting}

\subsection{Dobrushin uniqueness as the collapse-direction hypothesis}

The classical Dobrushin uniform-influence condition is packaged as the
predicate \lid{PaperUniformSmallTotalInfluence} on the MLN spec
in \codepath{Mettapedia.Logic.MarkovLogicInfiniteUniqueness}.
The collapse direction of the DLR-level dichotomy is the theorem
\lid{dlrQueryDetermined\_of\_uniform}:
\begin{lstlisting}
theorem dlrQueryDetermined_of_uniform
    (M : ClassicalInfiniteGroundMLNSpec Atom ClauseId)
    (hM : M.PaperUniformSmallTotalInfluence)
    (q : ConstraintQuery Atom) :
    dlrQueryDetermined M q
\end{lstlisting}
\emph{Under the Dobrushin uniform small-influence hypothesis, every finite
query is DLR-determined.}
Combining with the dichotomy gives the precision of the MLN query
envelope:
\begin{lstlisting}
theorem infiniteMLN_queryEnvelope_precise_of_uniform
    (M : ClassicalInfiniteGroundMLNSpec Atom ClauseId)
    [Nonempty (DLRCompletion M)]
    (hM : M.PaperUniformSmallTotalInfluence)
    (q : ConstraintQuery Atom) :
    infiniteMLNLowerQueryEnvelope M q
      = infiniteMLNUpperQueryEnvelope M q
\end{lstlisting}
\emph{Under the same hypothesis, and assuming the DLR completion family is
inhabited, the PLN $\STV{s}{c}$ for every finite query is precise.}

\subsection{Strict width as the phase-transition canary}

The opposite direction is the theorem
\lid{infiniteMLN\_queryEnvelope\_nontrivial\_of\_disagreement}
and the lift through
\lid{infiniteMLN\_queryEnvelope\_nontrivial\_of\_strictWidth}:
\begin{lstlisting}
theorem infiniteMLN_queryEnvelope_nontrivial_of_strictWidth
    (M : ClassicalInfiniteGroundMLNSpec Atom ClauseId)
    (q : ConstraintQuery Atom)
    (hBddBelow : BddBelow (Set.range (dlrCompletionQueryProb M q)))
    (hBddAbove : BddAbove (Set.range (dlrCompletionQueryProb M q)))
    (hWidth : dlrQueryHasStrictWidth M q) :
    infiniteMLNLowerQueryEnvelope M q
      < infiniteMLNUpperQueryEnvelope M q
\end{lstlisting}
\emph{If two DLR completions disagree on a finite query, the PLN
$\STV{s}{c}$ for that query is genuinely an interval.}

\subsection{Dobrushin collapse and strict-width witnesses}

Combining the two proved directions gives a one-sided collapse theorem
and a separate strict-width witness theorem. The Dobrushin condition
\lid{PaperUniformSmallTotalInfluence} is a formally sufficient hypothesis
for query determinacy and precise PLN/Walley envelopes. Separately, if
two DLR completions are exhibited that disagree on a query, then that
query has strict width. The Lean development does not prove that failure
of the Dobrushin hypothesis entails non-uniqueness or strict width.

\section{Worked examples already in the formalization}\label{sec:examples}

The codebase now ships both finite toy witnesses and a concrete infinite
DLR/PLN contrast crown. The finite two-spin examples remain useful sanity
witnesses for strict credal width.

\subsection{Two-spin magnet}

In ProjectiveCredal.lean around lines~1064--1101, the type
\lid{TwoSpin}, defined as $\mathrm{Bool} \times \mathrm{Bool}$, carries the two-spin
magnet credal set
\begin{lstlisting}
def twoSpinMagnetCredalSet : CredalPrevisionSet TwoSpin :=
  { PrecisePrevision.dirac twoSpinAllDown,
    PrecisePrevision.dirac twoSpinAllUp }
\end{lstlisting}
The credal set contains two Dirac measures: all spins down, and all
spins up. The first-spin-up gamble
\begin{lstlisting}
def twoSpinFirstUpGamble : Gamble TwoSpin :=
  fun s => if s.fst then 1 else 0
\end{lstlisting}
is rewarded by the all-up Dirac and unrewarded by the all-down Dirac;
the disagreement gives the strict width witness. The packaged theorem
\lid{twoSpinMagnetEnvelope\_nontrivial} (line~1101) closes:
\begin{lstlisting}
theorem twoSpinMagnetEnvelope_nontrivial :
    lowerEnvelope twoSpinMagnetCredalSet twoSpinFirstUpGamble
      < upperEnvelope twoSpinMagnetCredalSet twoSpinFirstUpGamble
\end{lstlisting}
PLN-side: on this two-Dirac credal, the magnetization query has lower
envelope $0$, upper envelope $1$, and confidence at its floor.

\subsection{Zero-temperature aligned magnet}

An incremental sibling
\lid{twoSpinZeroTemperatureMagnetEnvelope\_nontrivial} adds the
zero-temperature Gibbs constraint (spins must be aligned), and shows
the envelope is still strictly non-trivial: any measure supported on
the aligned-spin pair $\{(\downarrow,\downarrow), (\uparrow,\uparrow)\}$
still admits both extremes. The structural constraint by itself does
not close down credal width; the Dobrushin condition has to be checked
\emph{quantitatively} on the local-influence row-sums.

\subsection{Concrete infinite crown: reinforced line versus zero-weight grid}

The theorem
\lid{reinforcedLineGeometric\_strictWidth\_and\_zeroWeightGrid\_collapse}
packages the first concrete DLR/PLN contrast. The strict-width side uses
\lid{reinforcedLineGeometric\_originSpinUp\_plnStrictInterval}: a
geometrically reinforced half-line has positive PLN strict-width at the
origin-spin query. The collapse side uses
\lid{gridZeroField\_originSpinUp\_queryEnvelope\_precise\_of\_dobrushin}:
the zero-weight zero-field grid lies in the Dobrushin collapse regime for
the analogous origin-spin query. This is an infinite canary, but it is not
yet the full two-dimensional low-temperature Ising plus/minus
phase-separation theorem.

\section{Three kinds of degree of freedom: the synthesis}\label{sec:synthesis}

We now name the three formally-locatable freedom slots, each with the
predicate that witnesses it.

\paragraph{(A) Specification-level: \lid{localCredal} choice.}
A wider local credal at any cylinder window produces a wider projective
limit. The PLN-side question ``does my $\STV{s}{c}$ on query $X_q$ have
genuine width?'' reduces to ``does \lid{hasStrictGlobalWidth}
\lid{spec} on $X_q$ hold?''. The modeller controls this slot by varying
\lid{localCredal}{} $i$. The slot is finite or infinite depending on
whether $\Omega$ is.

\paragraph{(B) Completion-level: extreme-point choice.}
Given a credal set $C$, the actual world chooses one $P \in C$ (the
``true'' precise prevision). This freedom is aleatoric. It is formally
invisible to $\STV{s}{c}$ because $\STV{s}{c}$ is a property of $C$, not
of any specific $P$. PLN inference rules do not see it; only the
verification of a query against actual data sees it.

\paragraph{(C) Phase-transition: DLR non-uniqueness.}
In infinite Markov-logic specs, the projective limit is the set of DLR
completions, which may itself be non-singleton. The slot is opened
exactly when \lid{dlrQueryHasStrictWidth} on $(M, q)$ holds for the
query $q$ of interest, and closed when \lid{dlrQueryDetermined} on $(M, q)$
holds. The Dobrushin condition \lid{PaperUniformSmallTotalInfluence}
is a sufficient condition for the closed case.

\begin{observation}
Slot (A) is the freedom the PLN modeller intends to give. Slot (B) is
the freedom the world has within the modeller's spec. Slot (C) is the
freedom that the infinite structure of the spec itself produces, without
any modeller intent and without any world-side aleatoric residual.
\end{observation}

The three slots are independent. A modeller can pick (A) tight, the
world can pick (B) freely, and the spec can still have (C) freedom from
a phase transition. The PLN $\STV{s}{c}$ sees (A) and (C); it does not
see (B).

\section{What is open}\label{sec:open}

The formal characterization above is honest about which steps are
already-proved theorems and which are open. Since the first draft of this
section, the generic compact/FIP and Walley projective-natural-extension
waist has landed. The remaining work is now narrower:

\begin{itemize}
\item \textbf{Local finite-window coherence feeding FIP.}
      The generic compact/FIP waist is present:
      \lid{ProjectiveLocalLowerPrevisionSpec.finiteWindowCompatibleInCarrier}
      feeds
      \lid{hasCompatibleCompletion\_of\_finiteWindowCompatibleInCarrier}.
      What remains is an honest, non-circular local-coherence predicate
      that implies finite-window compatibility without assuming a hidden
      global compatible precise prevision.

\item \textbf{Polished DLR compact-projective packaging.}
      The DLR determined/strict-width bridge and the concrete
      reinforced-line crown exist. The remaining packaging work is to
      expose the DLR layer as a clean public specialization of the generic
      projective/Walley interface, with theorem names aligned to the
      confidence-formula API.

\item \textbf{Imprecise de~Finetti mixing-family package.}
      The finite-prefix Bernoulli-mixture payoff theorems exist, including
      agreement $\Rightarrow$ width-complement $1$ and disagreement
      $\Rightarrow$ width-complement $<1$. What remains is a clean public
      mixing-family object and, later, full process-law synthesis beyond
      finite-prefix adapters.

\item \textbf{True two-dimensional low-temperature Ising separation.}
      The reinforced-line strict-width crown is a concrete infinite canary.
      It is not yet the classical 2D low-temperature Ising plus/minus
      theorem. That theorem would require constructing two distinct
      infinite-volume DLR completions satisfying the same local conditional
      specification and proving that they disagree on an explicit
      finite-cylinder query.

\item \textbf{Paper/build metadata refresh.}
      Line counts and theorem references should be generated from the exact
      commit or avoided in favor of module and theorem names.
\end{itemize}

These are not blockers for the dichotomy we have proved; they are
the next theorems in the same vein. We name them as open so a reader
of this paper is not misled into thinking the infinite half of the
characterization is closed.

\section{Related work}\label{sec:related}

\paragraph{Walley.}
The credal-set lower-prevision framework is Walley (1991). Our
\lid{LowerPrevision} structure is the Walley structure with
constant-preservation, positive homogeneity, and superadditivity, and
our \lid{lowerEnvelopePrevision\_isCoherent\_of\_finite} is the
classical Walley coherence theorem for finite envelopes.

\paragraph{de~Finetti.}
The de~Finetti exchangeability theorem connects exchangeable distributions
on infinite sequences to mixtures over IID distributions. The codebase has
finite-prefix imprecise de~Finetti payoff theorems: mixture agreement
forces width-complement $1$, and mixture disagreement forces
width-complement $<1$ on the induced finite-prefix envelope. It also has a
projective specialization adapter. What remains open is a polished public
mixing-family object and a full process-law synthesis beyond the
finite-prefix surface.

\paragraph{Dobrushin--Lanford--Ruelle.}
The DLR framework characterises infinite Gibbs measures by their
finite-volume conditional distributions. The Dobrushin uniqueness
condition (uniform smallness of total influence) implies a unique DLR
completion. Singla and Domingos have a Markov-logic specific
uniformity-of-influence framing. Our \lid{PaperUniformSmallTotalInfluence}
is the codified hypothesis.

\paragraph{PLN literature.}
The two-component truth value $\STV{s}{c}$ is canonical in the PLN
literature (Goertzel et al.\ \emph{Probabilistic Logic Networks}). The
sibling papers \xiPLN\ (\lpath{papers/xiPLN.tex}) and the PLN Review
(\lpath{papers/pln-review.tex}) treat the evidence-algebra and
formal-rule-coherence axes. The WM-PLN book
(\lpath{papers/wm-pln-book.tex}) has the informal predecessor of the
degrees-of-freedom characterization in its \S 7.4
(``Degrees of freedom versus forcing'').

\paragraph{Markov logic networks.}
Richardson and Domingos (Machine Learning, 2006) introduced MLNs as
finite log-linear models on first-order logic. The infinite
specialisation in
\codepath{Mettapedia.Logic.MarkovLogicInfinite*}
is a recent codex landing (cited inline above).

\paragraph{Imprecise probability and second-order Bayesianity.}
The PLN tradition of treating $c$ as ``confidence-in-strength'' has a
natural Bayesian second-order reading (a Beta distribution over the
true probability), which the file
\codepath{Mettapedia.Logic.PLNDistributional} formalises. The credal
reading we develop here is the Walley-side complement; in the
finite case the two readings overlap on the singleton case and
diverge on the strict-width case.

\section{Conclusion}\label{sec:conclusion}

PLN truth values $\STV{s}{c}$ are summary statistics of a credal set,
and the question of where their information lives reduces, formally,
to a dichotomy on the underlying credal set. The codebase landings in
\codepath{Mettapedia.ProbabilityTheory.ImpreciseProbability.ProjectiveCredal}
	(the \lid{credalSetDetermines} / \lid{credalSetHasStrictWidth}
	predicates and their collapse / width theorems) and in
	\codepath{Mettapedia.Logic.MarkovLogicInfiniteCredalBridge} (the
	DLR-level lift and the Dobrushin hypothesis as a sufficient collapse
	condition) make this dichotomy a checkable Lean theorem rather than
	informal commentary.

The three freedom slots --- specification, completion, phase-transition
--- partition the question along orthogonal axes. (A) is the modeller's
freedom. (B) is the world's freedom. (C) is the structural freedom
the spec produces by itself. PLN $\STV{s}{c}$ sees (A) and (C); it
does not see (B). Where it sees (C), the Dobrushin condition is the
formally sufficient hypothesis for $\STV{s}{c}$ to be precise.

The infinite half of the characterization has real open work, named in
Sec.~\ref{sec:open}. We have not pretended otherwise.

\bibliographystyle{plain}
\begin{thebibliography}{10}

\bibitem{walley91}
Peter Walley.
\newblock \emph{Statistical Reasoning with Imprecise Probabilities}.
\newblock Chapman \& Hall, 1991.

\bibitem{definetti37}
Bruno de Finetti.
\newblock La pr\'evision: ses lois logiques, ses sources subjectives.
\newblock \emph{Annales de l'Institut Henri Poincar\'e}, 7:1--68, 1937.

\bibitem{dobrushin68}
R.\ L.\ Dobrushin.
\newblock Description of a random field by means of conditional
        probabilities and conditions of its regularity.
\newblock \emph{Theory of Probability and its Applications},
        13(2):197--224, 1968.

\bibitem{lanford-ruelle69}
O.\ Lanford and D.\ Ruelle.
\newblock Observables at infinity and states with short-range correlations
        in statistical mechanics.
\newblock \emph{Communications in Mathematical Physics}, 13:194--215, 1969.

\bibitem{richardson-domingos06}
Matthew Richardson and Pedro Domingos.
\newblock Markov logic networks.
\newblock \emph{Machine Learning}, 62:107--136, 2006.

\bibitem{singla-domingos05}
Parag Singla and Pedro Domingos.
\newblock Discriminative training of Markov logic networks.
\newblock In \emph{AAAI 2005}.

\bibitem{goertzel-pln}
Ben Goertzel, Matthew Ikl\'e, Izabela Goertzel, and Ari Heljakka.
\newblock \emph{Probabilistic Logic Networks: A Comprehensive Framework
        for Uncertain Inference}.
\newblock Springer, 2008.

\bibitem{xipln}
Mettapedia Working Notes.
\newblock $\xi$PLN: a correct and complete foundation for probabilistic
        logic networks.
\newblock In this directory: \lpath{papers/xiPLN.tex}.

\bibitem{wmpln-book}
Mettapedia Working Notes.
\newblock The WM-PLN book.
\newblock In this directory: \lpath{papers/wm-pln-book.tex}.

\bibitem{pln-review}
Mettapedia Working Notes.
\newblock PLN review and corrections.
\newblock In this directory: \lpath{papers/pln-review.tex}.

\end{thebibliography}

\end{document}
